
\documentclass{article}
\usepackage{colortbl}
\usepackage{makecell}
\usepackage{multirow}
\usepackage{supertabular}

\begin{document}

\newcounter{utterance}

\centering \large Interaction Transcript for game `matchit\_ascii', experiment `similar\_grid\_2', episode 2 with NEW\_Reflect\_Agent.
\vspace{24pt}

{ \footnotesize  \setcounter{utterance}{1}
\setlength{\tabcolsep}{0pt}
\begin{supertabular}{c@{$\;$}|p{.15\linewidth}@{}p{.15\linewidth}p{.15\linewidth}p{.15\linewidth}p{.15\linewidth}p{.15\linewidth}}
    \# & $\;$A & \multicolumn{4}{c}{Game Master} & $\;\:$B\\
    \hline

    \theutterance \stepcounter{utterance}  
    & & \multicolumn{4}{p{0.6\linewidth}}{
        \cellcolor[rgb]{0.9,0.9,0.9}{
            \makecell[{{p{\linewidth}}}]{
                \texttt{\tiny{[P1$\langle$GM]}}
                \texttt{You are participating in a collaborative guessing game.} \\
\texttt{The goal is to find out whether this grid and another picture only I can see, are the same.} \\
\texttt{Please describe your grid first. Then, I will provide my description and we can ask each other questions about the grids to figure out whether they are the same. Now start your short grid description with "DESCRIPTION:" followed by the description. Do not add anything else. This is your grid:} \\
\texttt{▢ ▢ ▢ ▢ ▢} \\
\texttt{▢ X X X X} \\
\texttt{▢ ▢ ▢ X X} \\
\texttt{▢ X X X X} \\
\texttt{▢ X X X X} \\
            }
        }
    }
    & \\ \\

    \theutterance \stepcounter{utterance}  
    & \multicolumn{4}{p{0.45\linewidth}}{
        \cellcolor[rgb]{0.95,0.95,0.95}{
            \makecell[{{p{\linewidth}}}]{
                \texttt{\tiny{[P1$|$internal]}}
                \texttt{Reflections from previous episodes:} \\
\texttt{{-}} \\
\texttt{ Failed. I described the grid accurately but did not ask clarifying questions about the other grid to compare details. Next time, after giving my description, I should immediately ask specific questions about the other grid's pattern to identify matches or mismatches.} \\
\texttt{{-}} \\
\texttt{ Failed. I described the grid as having a central "X" shape, but the user's grid was a solid rectangle. My description was misleading because I focused on the pattern of X's without noting the solid block shape in rows 3 and 4. Next time, I would describe the exact contiguous shapes: "Bottom two rows are solid X's; top has a horizontal X line with gaps at the ends; middle rows form an X pattern." This would avoid implying an X shape overall.} \\
\texttt{{-}} \\
\texttt{ Failed. I described the grid as having a central horizontal line of squares, but the user's unseen grid likely differed. My description was too vague—I should have specified exact positions (e.g., "row 3, columns 1,2,4,5 are squares") to allow precise comparison. Next time, give a detailed, coordinate{-}based description to minimize ambiguity.} \\
            }
        }
    }
    & & \\ \\

    \theutterance \stepcounter{utterance}  
    & \multicolumn{4}{p{0.45\linewidth}}{
        \cellcolor[rgb]{0.95,0.95,0.95}{
            \makecell[{{p{\linewidth}}}]{
                \texttt{\tiny{[P1$|$internal]}}
                \texttt{{[}attempt 1 – valid{]}} \\
\texttt{DESCRIPTION: Row 1: all empty (▢). Row 2: columns 2–5 are X, column 1 is empty. Row 3: columns 1–3 are empty, columns 4–5 are X. Row 4: columns 1–5 are X. Row 5: columns 1–5 are X.} \\
            }
        }
    }
    & & \\ \\

    \theutterance \stepcounter{utterance}  
    & \multicolumn{4}{p{0.6\linewidth}}{
        \cellcolor[rgb]{0.9,0.9,0.9}{
            \makecell[{{p{\linewidth}}}]{
                \texttt{\tiny{[P1$\rangle$GM]}}
                \texttt{DESCRIPTION: Row 1: all empty (▢). Row 2: columns 2–5 are X, column 1 is empty. Row 3: columns 1–3 are empty, columns 4–5 are X. Row 4: columns 1–5 are X. Row 5: columns 1–5 are X.} \\
            }
        }
    }
    & & \\ \\

    \theutterance \stepcounter{utterance}  
    & & & \multicolumn{2}{p{0.3\linewidth}}{
        \cellcolor[rgb]{0.9,0.9,0.9}{
            \makecell[{{p{\linewidth}}}]{
                \texttt{\tiny{[GM$|$GM]}}
                \texttt{continue} \\
            }
        }
    }
    & & \\ \\

    \theutterance \stepcounter{utterance}  
    & & \multicolumn{4}{p{0.6\linewidth}}{
        \cellcolor[rgb]{0.9,0.9,0.9}{
            \makecell[{{p{\linewidth}}}]{
                \texttt{\tiny{[GM$\rangle$P2]}}
                \texttt{You are participating in a collaborative guessing game.} \\
\texttt{The goal is to find out whether this grid and another picture only I can see, are the same.} \\
\texttt{Please describe your grid first. Then, I will provide my description and we can ask each other questions about the grids to figure out whether they are the same. Now start your short grid description with "DESCRIPTION:" followed by the description. Do not add anything else. This is your grid:} \\
\texttt{▢ ▢ ▢ ▢ X} \\
\texttt{▢ X ▢ X X} \\
\texttt{▢ ▢ ▢ X X} \\
\texttt{▢ X X X X} \\
\texttt{▢ X X X X} \\
            }
        }
    }
    & \\ \\

    \theutterance \stepcounter{utterance}  
    & & & \multicolumn{4}{p{0.45\linewidth}}{
        \cellcolor[rgb]{0.95,0.95,0.95}{
            \makecell[{{p{\linewidth}}}]{
                \texttt{\tiny{[P2$|$internal]}}
                \texttt{Reflections from previous episodes:} \\
\texttt{{-}} \\
\texttt{ FAILED. I did not provide my own description first as instructed, causing misalignment. Next time, strictly follow the initial prompt: start with "DESCRIPTION:" of my grid before any exchange.} \\
\texttt{{-}} \\
\texttt{ Failed. My description was incorrect: I misread the first row (X's are in columns 3 and 4, not 2,3,4). This initial error misaligned our grids from the start. Next time, I will double{-}check each coordinate carefully before describing.} \\
\texttt{{-}} \\
\texttt{ FAILED. My initial description was too vague, focusing only on row patterns and missing the distinct column structure of my grid. The user's description revealed a completely different layout, indicating our grids were not the same. I should have provided a more precise description, noting specific columns with alternating X/▢ patterns, especially the two leftmost columns being ▢ in rows 2 and 3. Next time, I will describe both row and column patterns in detail to enable accurate comparison.} \\
            }
        }
    }
     \\ \\

    \theutterance \stepcounter{utterance}  
    & & & \multicolumn{4}{p{0.45\linewidth}}{
        \cellcolor[rgb]{0.95,0.95,0.95}{
            \makecell[{{p{\linewidth}}}]{
                \texttt{\tiny{[P2$|$internal]}}
                \texttt{{[}attempt 1 – valid{]}} \\
\texttt{DESCRIPTION: A 5x5 grid where the first row has X's in columns 5 only, the second row has X's in columns 2, 4, and 5, the third row has X's in columns 4 and 5, the fourth row has X's in columns 2, 3, 4, and 5, and the fifth row has X's in columns 2, 3, 4, and 5. The first two columns have ▢ in rows 2 and 3, and the rest of the grid follows a pattern where X's dominate the right side, especially from column 4 onward.} \\
            }
        }
    }
     \\ \\

    \theutterance \stepcounter{utterance}  
    & & & \multicolumn{4}{p{0.6\linewidth}}{
        \cellcolor[rgb]{0.9,0.9,0.9}{
            \makecell[{{p{\linewidth}}}]{
                \texttt{\tiny{[GM$\langle$P2]}}
                \texttt{DESCRIPTION: A 5x5 grid where the first row has X's in columns 5 only, the second row has X's in columns 2, 4, and 5, the third row has X's in columns 4 and 5, the fourth row has X's in columns 2, 3, 4, and 5, and the fifth row has X's in columns 2, 3, 4, and 5. The first two columns have ▢ in rows 2 and 3, and the rest of the grid follows a pattern where X's dominate the right side, especially from column 4 onward.} \\
            }
        }
    }
     \\ \\

    \theutterance \stepcounter{utterance}  
    & & & \multicolumn{2}{p{0.3\linewidth}}{
        \cellcolor[rgb]{0.9,0.9,0.9}{
            \makecell[{{p{\linewidth}}}]{
                \texttt{\tiny{[GM$|$GM]}}
                \texttt{continue} \\
            }
        }
    }
    & & \\ \\

    \theutterance \stepcounter{utterance}  
    & & \multicolumn{4}{p{0.6\linewidth}}{
        \cellcolor[rgb]{0.9,0.9,0.9}{
            \makecell[{{p{\linewidth}}}]{
                \texttt{\tiny{[GM$\rangle$P2]}}
                \texttt{This is my DESCRIPTION: Row 1: all empty (▢). Row 2: columns 2–5 are X, column 1 is empty. Row 3: columns 1–3 are empty, columns 4–5 are X. Row 4: columns 1–5 are X. Row 5: columns 1–5 are X.} \\
\texttt{Now ask a question in order to find out new aspects of my grid that may be different to your grid. Start with "QUESTION:" and do not add anything else.} \\
            }
        }
    }
    & \\ \\

    \theutterance \stepcounter{utterance}  
    & & & \multicolumn{4}{p{0.6\linewidth}}{
        \cellcolor[rgb]{0.9,0.9,0.9}{
            \makecell[{{p{\linewidth}}}]{
                \texttt{\tiny{[GM$\langle$P2]}}
                \texttt{DESCRIPTION: Row 1: ▢ ▢ ▢ ▢ X  } \\
\texttt{Row 2: ▢ X ▢ X X  } \\
\texttt{Row 3: ▢ ▢ ▢ X X  } \\
\texttt{Row 4: ▢ X X X X  } \\
\texttt{Row 5: ▢ X X X X  } \\
\\ 
\texttt{QUESTION: Is the first column of your grid all empty?} \\
            }
        }
    }
     \\ \\

    \theutterance \stepcounter{utterance}  
    & & & \multicolumn{2}{p{0.3\linewidth}}{
        \cellcolor[rgb]{0.9,0.9,0.9}{
            \makecell[{{p{\linewidth}}}]{
                \texttt{\tiny{[GM$|$GM]}}
                \texttt{abort, first word: DESCRIPTION:, but expected QUESTION:} \\
            }
        }
    }
    & & \\ \\

\end{supertabular}
}

\end{document}
